\documentclass[11pt,a4paper]{article}
\usepackage[T1]{fontenc}
\usepackage[utf8]{inputenc}
\usepackage{lmodern}
\usepackage[margin=25mm]{geometry}
\usepackage{amsmath,amssymb,booktabs,tabularx,xurl,hyperref}
\hypersetup{colorlinks=true,linkcolor=blue!50!black,urlcolor=blue!50!black,
  pdftitle={Guarded Consumers: Projection, Evaluation and Ordered Effects},
  pdfauthor={Mettapedia Working Notes}}
\setlength{\emergencystretch}{3em}
\setlength{\parindent}{0pt}
\setlength{\parskip}{5pt}
\newcommand{\code}[1]{\nolinkurl{#1}}
\title{Guarded Consumers\\
  \large Projection, Evaluation and Ordered Effects}
\author{Mettapedia Working Notes}
\date{31 October 2026}
\begin{document}
\maketitle

\begin{abstract}
Six Lean modules extend the stored-ID and named-shape-store theory with
contracts for guard compilation, binding projection and effect batching.
A short-circuit guard compiler preserves results and executed read certificates.
A separate bridge compares it with eager Boolean relations only for pure,
complete, singleton-Boolean operands. Projecting bindings retains guards,
effects and answer templates together, preserving complete execution even on
partial failure. Ordered insert-on-absence plans preserve authored effects;
a narrower commutative set fragment admits multiplicity summaries with exact
branch-sensitive answer counts. These are logical permission laws, not native
compiler extraction or runtime performance guarantees.
\end{abstract}

\section{Two evaluation regimes must not be confused}

The guard language has literals, resolved query atoms, negation, conjunction
and disjunction. A query returns one of a Boolean answer, a fault, or an
unavailable realization. The last case denotes an optimization boundary, not
source-level falsity or a new source error. A source spelling does not acquire
a Boolean meaning merely because a compiler recognizes its name.

\code{GuardConjunctionCompiler} independently interprets the short-circuit
language and compiles it to the shared \code{ReadOnlyQuery.Program} interface.
The proof preserves the answer and complete ordered read certificate.
Constant-controlled branches and double negations normalize without changing
that observation. A validated executed-read certificate permits reuse in a
new world, but only for the same admitted program and authority context.
Changes of definitions, profile or primitive interpretation are not licensed
by a binding-only key.

The inspected PeTTa Boolean-relation path evaluates its argument tuple before
using a truth table. It must therefore not be identified with this
short-circuit language without another proof. For a pure finite-answer source,
an eager conjunction is the Cartesian combination
\[
 E(g\mathbin{\wedge}h)
   = [\,a\mathbin{\wedge}b\mid a\in E(g),\ b\in E(h)\,],
\]
where occurrences and order are retained. \code{GuardEvaluationRegimes}
proves equality of eager and short-circuit \emph{answers} when every atom
produces exactly one Boolean. This bridge does not preserve all reads: a
false left operand can make the right query unnecessary.

Three controls delimit the qualification:
\begin{itemize}
\item an eager false-and-empty expression has no answers, not one false answer;
\item a false left answer and two right answers produce two false occurrences;
\item a fault or effect on the right cannot be erased just by Boolean algebra.
\end{itemize}
The first two are kernel-checked list-semantic counterexamples. A small runtime
control also observed right-operand effects, zero answers for an empty right
operand, and duplicate false answers on the inspected saved build. General
source extraction must additionally establish purity, termination, absence of
binding refinements and the active primitive/profile interpretation.

\section{Projection retains the complete consumer}

\code{GuardedConsumerProjection} represents literal or variable operands,
guard calls, ordered effects, and finite answer templates. Its support is
computed structurally:
\[
 \begin{split}
 S(\operatorname{answer}(T)) &= S(T),\\
 S(\operatorname{effect}(e,a,p)) &= S(a)\cup S(p),\\
 S(\operatorname{branch}(g,p,q)) &= S(g)\cup S(p)\cup S(q).
 \end{split}
\]
The implementation keeps lists of coordinates; repeated coordinates are
allowed. A finite capture stores their binding values. Restoring arbitrary
fallback values elsewhere preserves the consumer's complete result:
answer occurrences, final world, ordered query reads, ordered effect calls
and completion status. An effect can leave a partially changed world and
then fault; that state is preserved too.

This closes the gap between ``unused by the answer'' and ``unused by the
consumer.'' A variable read only by a guard is still necessary. Likewise,
equal returned answer lists do not justify dropping an effect argument.
The static support is conservative rather than asserted minimal: a separately
certified specialization can eliminate unreachable branches first.

Projection is not memoization. Even a consumer with empty binding support
may observe a mutable world. A negative control changes a zero-argument
membership query while retaining the same binding capture, changing the
answer. Reuse across worlds needs read validation, and reuse across program
or profile changes needs the corresponding authority checks.

\section{Ordered plans for insert-on-absence guards}

For a guard key $k$, the fresh branch runs when $k$ is absent. The required
locality contract is: on successful completion it makes $k$ present and leaves
the other guard keys unchanged. Other world effects may be arbitrary and
noncommutative. On failure its actual partial state is retained.

\code{OrderedInsertGuardBatch.plan} walks the candidates without performing
world effects, recording fresh or present tokens against a private simulated
key set. It records every present token too, because the other branch can
produce an answer. \code{execute} then runs these tokens in the original
order. The theorem \code{planned_refines} equates it with independent sequential
guard evaluation, including failures, unavailable outcomes, answer
multiplicity, ordered effects and final state.

The planner's final key set is not bulk-published. If an early branch fails,
later tokens do not execute and their keys do not become visible. The proof
allows a failed branch to have made partial changes, rather than silently
rolling them back. Batches can be split at any candidate boundary under the
same sequential order; reordering arbitrary effects is not justified.

The key must encode all consumer information needed for these branch
functions. The projection theorem provides a finite binding capture, but
does not by itself justify treating two rows as the same guarded action.
Locality of guard-set changes and the complete effect contract remain
separate admission conditions.

\code{FirstRepresentativeGuardBatch} separates the guard key from the complete
consumer row. Rows with the same guard key may carry different effect
arguments. An ordered plan retains the first actual row for the fresh branch
and each later row for its present branch, and still refines sequential
execution. This allows effect variables outside the guard key when enumeration
and occurrence identity are preserved. It does not permit an arbitrary
representative: reversing two rows with the same key but different payloads
can change the result.

\section{A wider commutative region with exact multiplicity}

For set-valued insertions and a count-only observer, stronger permission is
available. \code{GuardedSetMultiplicity} models a state consisting of guard
keys, a set of facts and an answer count. A fresh key inserts a fixed set of
payloads; a present key changes neither set. Each branch has its own answer
count. Steps commute in this fragment, and permutation of candidates
preserves the full modeled state.

For multiplicity $m>0$, let $f(k)$ and $p(k)$ be the fresh-branch and
present-branch answer counts. The exact contribution is
\[
 \begin{cases}
    m\,p(k), & k\text{ initially present},\\
    f(k)+(m-1)\,p(k), & k\text{ initially absent}.
 \end{cases}
\]
Zero multiplicity changes nothing. The formula is proved against repeated
one-at-a-time execution. A reference checker validates that grouped counts
expand to exactly the source occurrence bag; checked groups then preserve
the final state. Missing a duplicate is rejected.

This permits an empty present branch: with $f(k)=1$, $p(k)=0$, one hundred
occurrences of a new key produce one answer, whereas an already held key
produces none. It also preserves multiple answers within the fresh branch.
It does not deduplicate branch-internal effects or answers.

This region excludes ordered answer streams, arbitrary stateful effects,
removal, primitive faults and intermediate observers. Those need the ordered
model or additional laws. Natural-number arithmetic is used; native overflow
and saturation policies are not specified by the theorem. The grouping checker
is a logical reference, not a claim that the production grouping algorithm has
been verified.

\section{Integration and qualification}

The six modules under \code{Mettapedia/Machines/} are
\code{GuardConjunctionCompiler}, \code{GuardEvaluationRegimes},
\code{GuardedConsumerProjection}, \code{OrderedInsertGuardBatch} and
\code{GuardedSetMultiplicity} and \code{FirstRepresentativeGuardBatch}.
\code{GuardOptimizationAudit} inventories
module-owned mathematical bodies and their transitive axioms. Standard
\code{propext}, \code{Classical.choice} and \code{Quot.sound} are permitted;
admissions and additional axioms are not. Compiler runtime counterparts are
listed separately from mathematical bodies.

The immediate native use is a specification for combining filters with
existence guards, retaining all consumer variables, and counting grouped
guarded insertions correctly. It supplies no new surface builtin and no native
optimization by itself. The frontend must establish source/profile authority,
pure scalar operands and stable captured bindings. Mutation validation and
publication must respect the selected observer, and implementation-level
resource/refusal behavior must remain correct.

These results extend the contracts described in \code{NamedShapeStores.tex}
and \code{StoredIdUnification.tex}. Restricted transducers and CTL fixed points,
canonical decision-diagram operations, higher derivatives and optimal cuts
remain separate parts of the wider theory.

\end{document}
